\section{Conclusão}

Este trabalho teve como objetivo o desenvolvimento e a implementação de um circuito capaz de calcular o quadrado de números naturais, utilizando a linguagem de descrição de hardware VHDL e a plataforma \textit{Quartus}, com testes finais realizados na placa DE10-Lite. A partir de uma abordagem sistemática e modular, foi possível projetar, simular e validar um circuito funcional, atendendo aos requisitos especificados no enunciado do projeto.

Durante a realização deste trabalho, diversas etapas foram cumpridas com êxito, desde a parte teórica até a validação prática. Na fase de simulação, foram realizados testes abrangentes que permitiram comprovar a consistência do circuito em condições ideais. Em seguida, a implementação na placa DE10-Lite confirmou a funcionalidade do projeto num ambiente mais realista, demonstrando a correspondência entre os resultados teóricos e práticos.

Como resultados finais, o circuito apresentou um desempenho consistente e eficiente, sendo capaz de calcular corretamente o quadrado de números naturais de 4 bits. A exibição dos resultados nos \textit{displays} de 7 segmentos confirmou a precisão e o funcionamento esperado do projeto, atingindo plenamente os objetivos estabelecidos.

Por fim, realçamos a modularidade presente no nosso circuito, podendo ser aprimorado, futura e facilmente, para a inclusão de números com maior quantidade de bits ou a implementação de outras operações matemáticas, por exemplo. Esta característica revela um bom conhecimento e aproveitamento, diante de inúmeras implementações possíveis para resolução deste trabalho, e por esse motivo, destaca a eficácia e singularidade da implementação do nosso circuito.